\subsection{Global extension: 38 earthquakes}\label{sec:res_global}
On the 38-event set, HiPhyBoost transfers better than PC-GBM in 29 of 38 earthquakes (Figure~\ref{fig:global}a). The median leave-one-event-out AUC is 0.795 against 0.763, and the mean difference is $+0.013$ (95\% bootstrap interval 0.006--0.020; Wilcoxon signed-rank $p=0.0006$, sign test $p=0.002$). Because all folds share most of their training events, these tests treat the folds as more independent than they are, and the $p$-values should be read as approximate. The gain is small but consistent. It holds for the 29 events that were not used to develop the method (21 of 29 better; $p=0.014$). Among the 16 new Pacific Rim events it is 12 of 16 ($p=0.07$), and among the 13 events outside the Pacific 9 of 13 ($p=0.15$). In this larger training pool HiPhyBoost is also better than PC-GBM in eight of the original nine events ($p=0.02$, against six of nine in the nine-event analysis), and Kaik\={o}ura keeps its high transfer skill (AUC 0.917 against 0.892). The worst transfers are the 2011 Tohoku-Oki earthquake (AUC 0.55--0.56), whose landslides are scattered over a very large, mostly weakly affected domain, and the 2013 Gansu earthquake, in largely loess terrain, where HiPhyBoost is worse than PC-GBM by 0.05. Without mapping, point forecasts of the landslide total remain poor for both models: the median error is a factor of 2.9 for HiPhyBoost and 4.0 for PC-GBM, and only 13 of 38 HiPhyBoost forecasts fall within a factor of two. This confirms, on four times as many events, that the level of landsliding cannot be predicted from terrain and shaking alone.

The fitted offsets range from $-4.2$ (Southern Italy 1980) to $+2.9$ (Cinchona), with $\tau=1.73$, close to the nine-event value. Their strongest correlate is the area of the study domain ($r=-0.53$, $p=0.001$; Figure~\ref{fig:global}b). For most events the domain is defined from the inventory, so this mainly reflects how much weakly affected terrain the mappers' domain includes. Offsets are also lower for point catalogues ($r=-0.42$, $p=0.009$), larger magnitudes ($r=-0.42$, $p=0.008$; magnitude and domain area are themselves correlated) and normal-faulting events ($r=-0.39$, $p=0.014$). Point catalogues and normal faulting partly coincide, because several of the catalogues are historical Italian records. Offsets are higher where a mapping extent was published ($r=+0.30$, $p=0.07$) and after wetter-than-normal months: precipitation over the strongly shaken area in the month before the earthquake, relative to its climatology, correlates with the offset at $r=+0.37$ ($p=0.02$; $r=+0.26$ when averaged over the landslide area). No individual correlation survives a Bonferroni correction for the 12 covariates examined.

Conditioning the zero-shot prior on forecast-time covariates sharpens the forecasts for unseen events (Figure~\ref{fig:global}c; Supplementary Table~S5). With magnitude, depth, mechanism and ShakeMap station count (F1), the median width of the 90\% interval for the landslide total falls from a factor of 364 to 132 (bootstrap 95\% interval 124--140) at unchanged coverage (33 of 38 events). The log predictive density of the held-out offsets improves by 0.18 nats per event (paired $p=0.04$), far above the improvement obtained with randomly permuted covariates (95th percentile 0.03; permutation $p=0.005$). Adding the antecedent-wetness covariates of the shaken area (F3) improves the prior further: width 107, coverage 34 of 38, gain 0.21 nats (paired $p=0.03$; permutation $p=0.005$). Climate alone (F2) does not help (gain $-0.05$), and computing it over the landslide area instead of the shaken area changes little (Supplementary Table~S5). With 38 events the coverage estimates are uncertain (Wilson 95\% interval about 0.73--0.94 for 33 of 38), so the differences in coverage between priors are not significant; the reduction in width is. Post-event covariates that describe the inventory and its domain (P2) narrow the interval further (width 120, coverage 36 of 38, gain 0.20). They cannot be used for forecasting, but they confirm that a substantial part of the offset reflects how and where the landslides were mapped. For Kaik\={o}ura, the forecast-time prior (F3) gives an interval of 0.05--4.2 times the observed total, against 0.03--9.4 for the unconditional prior, and both contain the observed total. The median point forecast of the total does not improve materially with any covariate set (error factor 2.8--3.6 against 2.9). Event covariates therefore currently help mainly to quantify how uncertain a zero-shot forecast is, not to predict the level of landsliding exactly.

\begin{figure*}[t]\centering
\includegraphics[width=\textwidth]{fig11_global.pdf}
\caption{Global extension to 38 earthquakes. (a) Leave-one-event-out AUC of HiPhyBoost against PC-GBM; points above the diagonal favour HiPhyBoost. (b) HiPhyBoost event offsets (fit to all events) against the area of the study domain, by inventory type. (c) Median width (bars, with bootstrap 95\% intervals over events) and coverage (line, with Wilson 95\% intervals; dotted: nominal 0.90) of the zero-shot 90\% interval for the landslide total. Climate covariates are averaged over the strongly shaken area. The bars compare the unconditional prior with priors conditioned on forecast-time covariates (blue: F1 shaking and ShakeMap, F2 climate, F3 both) and on post-event inventory and domain covariates (orange: P2).}\label{fig:global}
\end{figure*}
